\section{Analysis of Selected Heuristics}
\label{sec:selected}
\subsection{Why Do Final Solutions Converge?}

Across the five frozen EVO-F winners on the same 60 primary graphs,
construction orders differ from GR+P in 161 of 300 run--graph observations.
Only 61 have different pre-pruning vertex sets, 44 have different final
sets, and 23 have different final costs. Orders are compared as complete
sequences on the original labels; membership is compared as sets.
These are descriptive frequencies over dependent observations, not
300 independent graph samples.

The stages discard different kinds of variation. In 100 observations,
construction order changes without changing the constructed set.
Of the 61 different constructed sets, 17 converge to the same final set
after deterministic pruning. Another 21 observations retain different
final sets with equal cost. Thus, trajectory diversity exceeds
final-solution diversity, which in turn exceeds objective diversity.
Pruning explains one part of this convergence; selection during
construction and equal-cost alternatives explain the rest. A different
scoring expression can therefore produce many different choices without
improving the measured objective.

\begin{samepage}
\subsection{Why Does Rarity's Form Matter?}
\label{sec:rarity_form}

Seven of the 15 LLM winners are the greedy seed. Within EVO-F, runs 0
and 3 select that seed, while the other three select
\begin{equation}
\begin{aligned}
f_1 &= (g+0.4r)/q,\\
f_2 &= (g+0.55r)/q,\\
f_4 &= g/q+0.5r/q-0.2o.
\end{aligned}
\label{eq:selected_forms}
\end{equation}
\end{samepage}
The selected non-seed EVO-F programs therefore belong to a narrow
additive-rarity family, with an overlap penalty in run 4. This
characterizes the selected programs, not the full range of forms the
model could generate.

For the additive component, the bonus relative to coverage gain is
\begin{equation}
f(v,U)=\frac{g_v+\alpha r_v}{q_v},
\qquad
\frac{\alpha r_v}{g_v}
=\frac{\alpha r_v|U|}{k_v}.
\label{eq:rarity_scale}
\end{equation}
Coverage gain is normalized by the remaining demand $|U|$, whereas
rarity averages inverse neighborhood sizes over the $k_v$ newly covered
vertices. The relative bonus can therefore grow with $|U|$ or decrease
with $k_v$, even at the same average rarity. During construction, both
the uncovered set and the composition of each candidate's newly covered
neighborhood change. Graph size alone does not determine the bonus.

Multiplicative rarity has a different interpretation. Substituting the
feature definitions into Eq.~\eqref{eq:rarity} gives
\begin{equation}
\frac{g_vr_v}{q_v}
=
\frac{\bar c}{|U|}
\frac{1}{c_v}
\sum_{u\in N[v]\cap U}\frac{1}{|N[u]|}.
\label{eq:weighted_coverage}
\end{equation}
The leading factor is shared by all candidates at a step, so RGR+P
ranks vertices by total inverse-neighborhood weight of newly covered
demand divided by candidate cost. Standard greedy values each newly
covered vertex equally; RGR+P assigns more weight to demand vertices
with fewer possible covering neighbors. Additive rarity instead places
an average-rarity bonus alongside normalized coverage. It is not the
same weighted-coverage calculation.

\begin{table}[!htbp]
\centering\small
\caption{Interventions on all five frozen EVO-F winners. Mean excess
(\%) is relative to pruned greedy on training (30 graphs), primary and BA
(60 graphs each). No program is refit or reselected.}
\label{tab:ablation}
\begin{tabular}{lrrr}
\toprule
Intervention & Train & Primary & BA\\
\midrule
% Generated from saved JSON by scripts/reproduce_results.py.
Locked & -0.144 & +0.313 & +1.438 \\
$r=0$ & +0.209 & +0.200 & +0.049 \\
$o=0$ & -0.154 & +0.141 & +1.430 \\
No pruning & +6.922 & +7.558 & +10.642 \\
\bottomrule

\end{tabular}
\end{table}

\begin{samepage}
Table~\ref{tab:ablation} summarizes all five winners. For the three
non-seed winners, zeroing rarity changes training excess from
$-0.240\%$ to $+0.349\%$, primary excess from $+0.522\%$ to
$+0.334\%$, and BA excess from $+2.396\%$ to $+0.082\%$.\par
\end{samepage}

\begin{samepage}
The rarity term helps training fitness but its contribution does not
transfer to these held-out averages. Runs 1 and 2 reduce exactly to
greedy when rarity is zeroed; run 4 retains its overlap penalty.
Omitting pruning also worsens the five-winner averages, showing the
importance of the fixed scaffold.\par
\end{samepage}

RGR+P, by contrast, improves mean cost on BA and both external weight
families (Tables~\ref{tab:shift} and~\ref{tab:external}).
Together, the algebra and interventions make the functional form of
rarity weighting a plausible explanation for transfer differences.
They analyze the selected heuristics' behavior without establishing a
unique cause of the discovery policies' outcomes.

\subsection{Illustrative Decision Trajectories}

The following saved traces are the first qualifying primary improvement
and BA failure in lexicographic (run, instance) order. They illustrate
how scoring changes interact with construction and pruning.

\paragraph{A primary improvement.}
For primary instance 0014, run 1 uses $f_1$ from
Eq.~\eqref{eq:selected_forms}. At the third selection, vertices 15 and 35
both cover seven new vertices at cost two. Greedy selects 15 by ID;
the higher average rarity makes EVO-F select 35. EVO-F finishes at
cost 29, the certified optimum, versus greedy's 31. Vertex 35 is
itself pruned: its effect is on the subsequent construction trajectory,
rather than final membership. The strict improvement survives one of
ten saved vertex relabelings; the other nine tie greedy, showing the
example's sensitivity to ID-based tie-breaking.

\paragraph{A BA failure.}
For BA instance 0044, run 4 uses $f_4$.
It agrees with greedy for four selections, then chooses vertex 30
instead of vertex 19. Their new-coverage-to-cost ratios are 1 and 2,
respectively. The additive rarity term reverses this ranking even
with the overlap penalty; without rarity, the two candidates retain
greedy's ordering. Both trajectories later prune vertex 30.
Final costs are 34 for EVO-F and 33 for greedy, while RGR+P reaches
the certified optimum of 30. The first ranking reversal identifies a
departure that propagates through later choices; it does not isolate
the entire final loss to one decision. Complete traces and per-winner
summaries remain in the repository.
